% ============================================================
% COMANDOS PERSONALIZADOS - trab-disc-unesp
% ============================================================

% ----- CONTROLE DE ATIVAÇÃO DOS PRÉ-TEXTUAIS -----
\newcommand{\exibirResumo}{%
    \ifnum\ativarResumo=1
        % ============================================================
% RESUMO - trab-disc-unesp
% ============================================================

\begin{resumo}

% --- TEXTO DO RESUMO (12pt, ESPAÇAMENTO SIMPLES, SEM RECUO, JUSTIFICADO) ---
\setlength{\parindent}{0pt}
\setlength{\parskip}{0pt}
\singlespacing
\normalsize

Escreva aqui o resumo em português. Deve conter entre 150 e 500 palavras em um único parágrafo. O resumo deve apresentar de forma concisa os objetivos, a metodologia, os resultados e as conclusões do trabalho.

\vspace{0.5cm}

% --- PALAVRAS-CHAVE (12pt, NEGRITO PARA O TERMO "Palavras-chave:") ---
{\normalsize\bfseries Palavras-chave:} \normalsize
\palavrachaveum. \palavrachavedois. \palavrachavetres.%
\ifx\palavrachavequatro\empty\else
    \ \palavrachavequatro.%
\fi
\ifx\palavrachavecinco\empty\else
    \ \palavrachavecinco.%
\fi

\end{resumo}

% ============================================================
% FIM DO ARQUIVO resumo.tex
% ============================================================
    \fi
}

\newcommand{\exibirAbstract}{%
    \ifnum\ativarResumo=1
        \ifnum\ativarAbstract=1
            % ============================================================
% ABSTRACT - trab-disc-unesp
% ============================================================

\chapter*{ABSTRACT}

\begin{center}
    {\fontspec{Times New Roman}\fontsize{14pt}{16pt}\bfseries\selectfont%
        \tituloingles%
    \par}
    \vspace{0.3cm}
    \ifx\subtitulo\empty\else
        {\fontspec{Times New Roman}\fontsize{12pt}{14pt}\bfseries\selectfont%
            \subtitulo%
        \par}%
        \vspace{0.3cm}
    \fi
\end{center}

\vspace{0.5cm}

{\fontspec{Times New Roman}\fontsize{12pt}{14pt}\selectfont%
    Write the abstract in English here. It should contain between 150 and 500 words in a single paragraph.
\par}

\vspace{0.5cm}

{\fontspec{Times New Roman}\fontsize{12pt}{14pt}\bfseries\selectfont%
    Keywords: %
    \fontspec{Times New Roman}\fontsize{12pt}{14pt}\selectfont%
    \keywordum. \keyworddois. \keywordtres.%
    \ifx\keywordquatro\empty\else
        \ \keywordquatro.%
    \fi
    \ifx\keywordcinco\empty\else
        \ \keywordcinco.%
    \fi
\par}

\vspace{0.5cm}

{\fontspec{Times New Roman}\fontsize{12pt}{14pt}\bfseries\selectfont%
    Title: %
    \fontspec{Times New Roman}\fontsize{12pt}{14pt}\selectfont%
    \tituloingles%
\par}

% ============================================================
% FIM DO ARQUIVO abstract.tex
% ============================================================
        \fi
    \fi
}

% ----- COMANDO PARA EXIBIR SUMÁRIO (SEM NUMERAÇÃO) -----
\newcommand{\exibirSumario}{%
    \ifnum\ativarSumario=1
        \vspace*{0pt}
        \begin{center}
            {\fontsize{14pt}{16pt}\selectfont\bfseries SUMÁRIO}
        \end{center}
        \vspace{-3cm}
        \renewcommand{\contentsname}{}
        \thispagestyle{empty}  % PRIMEIRA VEZ
        \tableofcontents
        \thispagestyle{empty}  % SEGUNDA VEZ (FORÇA APÓS O \tableofcontents)
        \clearpage
    \fi
}

% ----- COMANDOS PARA PÓS-TEXTUAIS -----
\newcommand{\exibirApêndices}{%
    \ifnum\ativarApêndices=1
        % ============================================================
% APÊNDICES - trab-disc-unesp
% ============================================================

\chapter{APÊNDICE A - Título do Apêndice}
\label{apendice:a}

Texto do apêndice...

\chapter{APÊNDICE B - Título do Apêndice}
\label{apendice:b}

Texto do apêndice...

% ============================================================
% FIM DO ARQUIVO 02-apendices.tex
% ============================================================
    \fi
}

\newcommand{\exibirAnexos}{%
    \ifnum\ativarAnexos=1
        % ============================================================
% ANEXOS - trab-disc-unesp
% ============================================================

\chapter{ANEXO A - Título do Anexo}
\label{anexo:a}

Texto do anexo...

\chapter{ANEXO B - Título do Anexo}
\label{anexo:b}

Texto do anexo...

% ============================================================
% FIM DO ARQUIVO 01-anexos.tex
% ============================================================
    \fi
}

% ----- AMBIENTES PRÉ-TEXTUAIS -----
\newenvironment{resumo}{%
    \thispagestyle{empty}
    \vspace*{0pt}
    \begin{center}
        {\fontsize{14pt}{16pt}\selectfont\bfseries RESUMO}
    \end{center}
    \vspace{0.5cm}
}{%
    \par
    \clearpage
}

\newenvironment{abstract}{%
    \thispagestyle{empty}
    \vspace*{0pt}
    \begin{center}
        {\fontsize{14pt}{16pt}\selectfont\bfseries ABSTRACT}
    \end{center}
    \vspace{0.5cm}
}{%
    \par
    \clearpage
}

% ----- COMANDO PARA FORMATAR TÍTULO DO DOCENTE -----
\newcommand{\formatarDocente}[3]{%
    \ifthenelse{\equal{#2}{f}}{%
        \ifthenelse{\equal{#3}{1}}{Profa. Dra. }{Profa. }%
    }{%
        \ifthenelse{\equal{#3}{1}}{Prof. Dr. }{Prof. }%
    }%
    #1%
}

% ----- COMANDO PARA LISTAR DOCENTES NA CAPA -----
\newcounter{totalDocentes}
\newcommand{\listarDocentes}{%
    \setcounter{totalDocentes}{1}
    \ifx\docenteDoisNome\undefined\else
        \ifx\docenteDoisNome\empty\else
            \stepcounter{totalDocentes}
        \fi
    \fi
    \ifx\docenteTresNome\undefined\else
        \ifx\docenteTresNome\empty\else
            \stepcounter{totalDocentes}
        \fi
    \fi
    
    \ifnum\value{totalDocentes}=1
        \textbf{Docente:} \formatarDocente{\docenteUmNome}{\docenteUmGenero}{\docenteUmDoutor}
    \else
        \textbf{Docentes:} \formatarDocente{\docenteUmNome}{\docenteUmGenero}{\docenteUmDoutor}
    \fi
    
    \ifx\docenteDoisNome\undefined\else
        \ifx\docenteDoisNome\empty\else
            \par
            \formatarDocente{\docenteDoisNome}{\docenteDoisGenero}{\docenteDoisDoutor}
        \fi
    \fi
    
    \ifx\docenteTresNome\undefined\else
        \ifx\docenteTresNome\empty\else
            \par
            \formatarDocente{\docenteTresNome}{\docenteTresGenero}{\docenteTresDoutor}
        \fi
    \fi
}

% ----- COMANDO PARA LISTAR AUTORES NA CAPA -----
\newcounter{totalAutoresTemp}
\newcommand{\listarAutores}{%
    \setcounter{totalAutoresTemp}{1}
    \ifx\coautorUm\undefined\else\ifx\coautorUm\empty\else\stepcounter{totalAutoresTemp}\fi\fi
    \ifx\coautorDois\undefined\else\ifx\coautorDois\empty\else\stepcounter{totalAutoresTemp}\fi\fi
    \ifx\coautorTres\undefined\else\ifx\coautorTres\empty\else\stepcounter{totalAutoresTemp}\fi\fi
    \ifx\coautorQuatro\undefined\else\ifx\coautorQuatro\empty\else\stepcounter{totalAutoresTemp}\fi\fi
    \ifx\coautorCinco\undefined\else\ifx\coautorCinco\empty\else\stepcounter{totalAutoresTemp}\fi\fi
    \ifx\coautorSeis\undefined\else\ifx\coautorSeis\empty\else\stepcounter{totalAutoresTemp}\fi\fi
    \ifx\coautorSete\undefined\else\ifx\coautorSete\empty\else\stepcounter{totalAutoresTemp}\fi\fi
    \ifx\coautorOito\undefined\else\ifx\coautorOito\empty\else\stepcounter{totalAutoresTemp}\fi\fi
    \ifx\coautorNove\undefined\else\ifx\coautorNove\empty\else\stepcounter{totalAutoresTemp}\fi\fi
    
    \ifnum\value{totalAutoresTemp}=1
        \textbf{Nome:} \autorPrincipal{} \textbf{- RA:} \raPrincipal
    \else
        \textbf{Discentes:} \\
        \autorPrincipal{} \textbf{- RA:} \raPrincipal
        \ifx\coautorUm\undefined\else\ifx\coautorUm\empty\else\\ \coautorUm{} \textbf{- RA:} \raCoautorUm\fi\fi
        \ifx\coautorDois\undefined\else\ifx\coautorDois\empty\else\\ \coautorDois{} \textbf{- RA:} \raCoautorDois\fi\fi
        \ifx\coautorTres\undefined\else\ifx\coautorTres\empty\else\\ \coautorTres{} \textbf{- RA:} \raCoautorTres\fi\fi
        \ifx\coautorQuatro\undefined\else\ifx\coautorQuatro\empty\else\\ \coautorQuatro{} \textbf{- RA:} \raCoautorQuatro\fi\fi
        \ifx\coautorCinco\undefined\else\ifx\coautorCinco\empty\else\\ \coautorCinco{} \textbf{- RA:} \raCoautorCinco\fi\fi
        \ifx\coautorSeis\undefined\else\ifx\coautorSeis\empty\else\\ \coautorSeis{} \textbf{- RA:} \raCoautorSeis\fi\fi
        \ifx\coautorSete\undefined\else\ifx\coautorSete\empty\else\\ \coautorSete{} \textbf{- RA:} \raCoautorSete\fi\fi
        \ifx\coautorOito\undefined\else\ifx\coautorOito\empty\else\\ \coautorOito{} \textbf{- RA:} \raCoautorOito\fi\fi
        \ifx\coautorNove\undefined\else\ifx\coautorNove\empty\else\\ \coautorNove{} \textbf{- RA:} \raCoautorNove\fi\fi
    \fi
}

% ----- COMANDO PARA CONTAR AUTORES -----
\newcounter{totalAutores}
\newcommand{\contarAutores}{%
    \setcounter{totalAutores}{1}
    \ifx\coautorUm\undefined\else\ifx\coautorUm\empty\else\stepcounter{totalAutores}\fi\fi
    \ifx\coautorDois\undefined\else\ifx\coautorDois\empty\else\stepcounter{totalAutores}\fi\fi
    \ifx\coautorTres\undefined\else\ifx\coautorTres\empty\else\stepcounter{totalAutores}\fi\fi
    \ifx\coautorQuatro\undefined\else\ifx\coautorQuatro\empty\else\stepcounter{totalAutores}\fi\fi
    \ifx\coautorCinco\undefined\else\ifx\coautorCinco\empty\else\stepcounter{totalAutores}\fi\fi
    \ifx\coautorSeis\undefined\else\ifx\coautorSeis\empty\else\stepcounter{totalAutores}\fi\fi
    \ifx\coautorSete\undefined\else\ifx\coautorSete\empty\else\stepcounter{totalAutores}\fi\fi
    \ifx\coautorOito\undefined\else\ifx\coautorOito\empty\else\stepcounter{totalAutores}\fi\fi
    \ifx\coautorNove\undefined\else\ifx\coautorNove\empty\else\stepcounter{totalAutores}\fi\fi
}

% ============================================================
% FORMATAÇÃO DE CAPÍTULOS E SEÇÕES
% ============================================================
\titleformat{\chapter}[display]
    {\normalfont\fontsize{14pt}{16pt}\selectfont\bfseries\centering}
    {}
    {0pt}
    {\MakeUppercase}

\titleformat{\section}
    {\normalfont\normalsize\bfseries}
    {\thesection}
    {0.5em}
    {}

\titleformat{\subsection}
    {\normalfont\normalsize\bfseries}
    {\thesubsection}
    {0.5em}
    {}

\titleformat{\subsubsection}
    {\normalfont\normalsize\itshape}
    {\thesubsubsection}
    {0.5em}
    {}

\titlespacing*{\chapter}{0pt}{0pt}{1cm}

% ============================================================
% CONFIGURAÇÃO DO CABEÇALHO (NUMERAÇÃO SUPERIOR DIREITA)
% ============================================================

\pagestyle{fancy}
\fancyhf{}  % Limpa todos os cabeçalhos e rodapés

% --- NUMERAÇÃO NO CANTO SUPERIOR DIREITO ---
\fancyhead[R]{\thepage}

% --- REMOVE A LINHA HORIZONTAL DO CABEÇALHO ---
\renewcommand{\headrulewidth}{0pt}

% --- GARANTE QUE O ESTILO fancy SEJA USADO EM TODAS AS PÁGINAS ---
\renewcommand{\chaptermark}[1]{}
\renewcommand{\sectionmark}[1]{}

% --- ESTILO PARA PÁGINAS DE CAPÍTULO (com numeração) ---
\fancypagestyle{plain}{%
    \fancyhf{}%
    \fancyhead[R]{\thepage}%
    \renewcommand{\headrulewidth}{0pt}%
}

% ----- COMANDOS PARA FIGURAS -----
\newcommand{\figura}[4]{%
    \begin{figure}[H]
        \centering
        \includegraphics[width=#1]{#2}
        \caption{#3}
        \label{fig:#4}
    \end{figure}
}

% ----- COMANDOS PARA TABELAS -----
\newcommand{\tabela}[4]{%
    \begin{table}[H]
        \centering
        \caption{#1}
        \label{tab:#2}
        #3
        \fonte{#4}
    \end{table}
}

\newcommand{\fonte}[1]{%
    \vspace{0.2cm}
    \par\small\textbf{Fonte:} #1
}

\newcommand{\nota}[1]{%
    \vspace{0.2cm}
    \par\small\textbf{Nota:} #1
}

\newcommand{\destaque}[1]{\textbf{#1}}
\newcommand{\citacao}[1]{\textit{#1}}

% ============================================================
% FIM DO ARQUIVO comandos.tex
% ============================================================